\documentclass[11pt]{article}
\usepackage{amsmath,amssymb,amsthm,mathtools}
\usepackage[T1]{fontenc}
\usepackage{lmodern}
\usepackage[margin=1in]{geometry}
\usepackage{enumitem}
\usepackage{booktabs}
\usepackage{array}
\usepackage{longtable}
\usepackage{microtype}
\usepackage{xcolor}
\usepackage{hyperref}
\usepackage{xurl}

\hypersetup{colorlinks=true,linkcolor=blue,citecolor=blue,urlcolor=blue}
\urlstyle{same}

\newtheorem{theorem}{Theorem}[section]
\newtheorem{proposition}[theorem]{Proposition}
\newtheorem{remark}[theorem]{Remark}
\newtheorem{definition}[theorem]{Definition}
\newtheorem{question}[theorem]{Open problem}

\newcommand{\Z}{\mathbb Z}
\newcommand{\wdone}{\widehat{\mathfrak N}_m(0)}
\newcommand{\wdthree}{\widehat{\mathfrak N}_m(1)}
\newcommand{\Pm}{P_m}
\newcommand{\SelStar}{\mathrm{Sel}^{\star}}
\newcommand{\status}[1]{\textsc{#1}}

\title{Companion research note for the residual assembly program\\
after the nonresonant closure}
\author{}
\date{March 25, 2026}

\begin{document}
\maketitle

\begin{abstract}
This note is meant to be read alongside the redefined theorem-order manuscript
for odd-$m$, $d=5$.
Its job is not to replace the paper, but to make the current proof status easy
for another researcher to follow.
The note records four things in one place:
(1) the manuscript redefinition after the closure of the small odd and
nonresonant residual packets;
(2) the exact reduced-return laws that make the nonresonant closure believable
and reusable;
(3) the current resonant reduction tree, including the deterministic reductions
from the checked $51$- and $69$-type branches to canonical pure color-$1$
frontiers; and
(4) the obstruction ladder that explains why the remaining resonant theorem is
now concentrated on a double-top breaker problem plus a separate
$B$-active/gate branch.
Throughout, the note distinguishes carefully between theorem-level results,
checked exact packets, and genuine frontiers.
\end{abstract}

\tableofcontents

\section{Purpose and reading protocol}

The redefined manuscript is now organized around the following theorem-order
spine:
\[
\text{front end}
\to
\text{odometer spine}
\to
\text{selector surgery}
\to
\text{generic upgrade}.
\]
What changed is the location of the honest boundary.
The paper no longer treats residual assembly as a single graph-side input for
all odd moduli.
Instead:
\begin{enumerate}[label=\textup{(\arabic*)},leftmargin=2.8em]
\item the small odd moduli $m=5,7,9$ are closed explicitly;
\item the nonresonant residual theorem for $3\nmid m$ is closed by the family
      \(\Pm\);
\item the only remaining external input is a \emph{resonant} residual theorem
      for odd $m\ge 15$ with $3\mid m$.
\end{enumerate}
Accordingly, the updated paper should be read as a theorem-order proof spine,
while the present note should be read as the companion packet that records the
new exact reduced-return formulas, the case table, and the remaining resonant
obstructions.

\begin{remark}[how the paper and note divide labor]
The paper keeps the load-bearing theorem chain short and legible.
This note records the exact reduced maps and branchwise status that justify the
promotion of the nonresonant closure and the narrowing of the resonant
frontier.
When a theorem in the paper is labeled by a short proof sketch or a companion
pointer, the detailed formula or checked branch table lives here.
\end{remark}

\section{Redefinition of the manuscript spine}

The most important editorial change is the redefinition of the old residual
assembly background input.
Table~\ref{tab:promotion-map} records the old role, the new role, and the
current status.

\begin{table}[h!]
\centering
\small
\begin{tabular}{p{0.24\linewidth}p{0.30\linewidth}p{0.17\linewidth}p{0.19\linewidth}}
\toprule
Old location / role & New role in the redefined manuscript & Status & Where the details now live \\
\midrule
Background 8.16: all odd residual assembly & Frontier theorem for the resonant range $m\ge 15$, $3\mid m$ only & \status{Frontier} & Updated paper, Section 8 \\
Graph-side burden after color-$3$/color-$4$ closure & Split into small odd closure, nonresonant theorem, and resonant frontier & mixed & Updated paper, Section 8 + this note \\
Final upgrade assumption (G3) & Resonant-only residual input & \status{Promoted} & Updated paper, Section 10 \\
Methodological picture & Resonant-only obstruction ladder: width-1, width-3, gate branch & \status{Promoted} & Updated paper, Section 9 + Sections 5--6 here \\
Residual proof packets for colors $0,1,2$ & Explicit theorem packets for small odd and nonresonant moduli & \status{Closed} & Section 3 here \\
\bottomrule
\end{tabular}
\caption{Promotion map after the nonresonant closure.}
\label{tab:promotion-map}
\end{table}

At the level of theorem dependence, the global statement is now:
\[
\text{main theorem}
=
\text{closed backbone}
+
\text{closed nonresonant residual theorem}
+
\text{open resonant residual theorem}.
\]
This is the central organizational fact behind the rewrite.

\section{Closed residual packets outside resonance}

\subsection{Small odd moduli}

The small odd moduli are closed by explicit packages:
\[
\begin{aligned}
&m=5:\quad \{(0,1,0,0),\ (0,2,1,0),\ (0,1,3,3)\},\\
&m=7:\quad \{(0,1,0,0),\ (0,2,1,0),\ (0,1,3,3),\ (1,2,3,3)\},\\
&m=9:\quad \{(0,1,4,4),\ (1,2,4,3)\}.
\end{aligned}
\]
These are theorem-level closures: exact finite replay on the full torus shows
that they preserve the frozen color-$3$/color-$4$ backbone and make the
residual colors $0,1,2$ Hamilton.

\subsection{The nonresonant family \texorpdfstring{$\Pm$}{Pm}}

For odd $m\ge 9$ let
\[
\Pm:=\{(1,2,1,3),\ (1,2,1,4),\ (0,1,(m-1)/2,4)\}.
\]
This family is now the closed theorem-side package for all odd moduli with
$3\nmid m$.
The three colors are handled by three separate reduced-return packets.

\paragraph{Color 1.}
The color-$1$ torus map is fully classified.
For $3\nmid m$ it is Hamilton.
For $3\mid m$ it has full-torus cycle lengths
\[
\frac{m^4(m+1)}3,\quad \frac{m^4(m-1)}3,\quad \frac{m^4(m-2)}3,
\quad \frac{m^4}3,\quad \frac{m^4}3.
\]
The proof descends through two exact clocks and an exact lower map
\(U_1\), which is generically a translation
\((a,c)\mapsto(a,c-4)\) and differs only on a finite defect frame.
The nonresonant branch closes because the row-$4$ first-return permutation on
that lower section is a single $m$-cycle with total return time $m^2$.

\paragraph{Color 2.}
Write
\[
U_2=T_2^m=h_2^{m^3}\big|_{\Sigma=0,\ x_3=0,\ x_4=0}
\]
in lower coordinates $(a,b)=(x_0,x_1)$.
The exact lower law is
\[
U_2(a,b)=
\begin{cases}
(a-4,b+2), & b\in\{0,3,m-3\},\\
(a-6,b+2), & b\in\{1,2\},\\
(a-5,b+2), & \text{otherwise.}
\end{cases}
\]
Since $b\mapsto b+2$ is a single $m$-cycle for odd $m$ and the total $a$-drift
around one full $b$-cycle is
\[
3(-4)+2(-6)+(m-5)(-5)=-5m+1\equiv 1\pmod m,
\]
one gets a single $m^2$-cycle on the lower section and therefore a Hamilton
full-torus map.

\paragraph{Color 0.}
The decisive object is the strip first-return
\[
F_0:S_0\to S_0,
\qquad
S_0=\{\Sigma=0,\ x_4=0,\ x_1=0\},
\]
with strip coordinates $(a,d)=(x_0,x_3)$.
The exact law is
\[
F_0(a,d)=
\begin{cases}
(0,m-1), & (a,d)=(m-1,m-1),\\
(a+4,0), & d=m-1,\ a\neq m-1,\\
(a+1,m-2), & d=m-2,\ a=4,\\
(a-4,m-1), & d=m-2,\ a\neq 4,\\
(a+1,d), & d\le m-3,\ a+d=2,\\
(a,d+1), & d\le m-3,\ a+d\neq 2.
\end{cases}
\]
The return-time function is
\[
\tau_0(a,d)=
\begin{cases}
(m+1)/2,& (a,d)=(m-1,m-1),\\
m+(m-1)/2,& (a,d)=(5,m-2),\\
m,& \text{otherwise.}
\end{cases}
\]
Every orbit reaches the top row $d=m-1$, the induced top-row first-return is
exactly $a\mapsto a+1$, and therefore $F_0$ is a single $m^2$-cycle on $S_0$.
Since the total time over the strip is $m^3$, the lifted section return is one
$m^3$-cycle, hence color $0$ is Hamilton.

\begin{proposition}[nonresonant residual theorem, packet form]
For every odd modulus $m\ge 11$ with $3\nmid m$, the family $\Pm$ leaves the
frozen color-$3$/color-$4$ backbone unchanged and makes colors $0,1,2$
Hamilton.
Consequently, together with the explicit closures at $m=5,7,9$, the entire
residual assembly program is closed outside the resonant range.
\end{proposition}

\section{Visibility, quotient reduction, and arithmetic markers}

The resonant search is no longer an unrestricted search over all row words.
A fixed-color quotient theorem cuts down the effective state space.

\subsection{Fixed-color quotients}

For a single row word
\[
\pi\in S_3=\{\mathrm{id},A=(12),B=(02),C=(01),\rho=BA,\rho^{-1}=AB\},
\]
the effect on a fixed color $c\in\{0,1,2\}$ depends only on
\(\pi^{-1}(c)\).
Thus each color sees only three quotient classes.
Writing the classes by the color they send into the observed color gives:
\[
Q_{0,0}=\{\mathrm{id},A\},\quad Q_{0,1}=\{C,\rho\},\quad Q_{0,2}=\{B,\rho^{-1}\},
\]
\[
Q_{1,0}=\{C,\rho^{-1}\},\quad Q_{1,1}=\{\mathrm{id},B\},\quad Q_{1,2}=\{A,\rho\},
\]
\[
Q_{2,0}=\{B,\rho\},\quad Q_{2,1}=\{A,\rho^{-1}\},\quad Q_{2,2}=\{\mathrm{id},C\}.
\]
This has two immediate consequences.

\begin{enumerate}[label=\textup{(\arabic*)},leftmargin=2.8em]
\item Same-row same-class changes are fake freedom for the observed color.
\item The three basic transpositions have channel visibility laws:
      \(A\) is invisible to color $0$, \(B\) is invisible to color $1$,
      and \(C\) is invisible to color $2$.
\end{enumerate}

In particular, no nontrivial single-row word can change exactly one color while
leaving the other two fixed.
That is why the resonant repairs are inherently staged: one channel fixes the
current bad color while collateral damage is passed to a second channel.

\subsection{Arithmetic case table}

The current resonant search is best organized by
\((m\bmod 24,\ 5\mid m)\).
Table~\ref{tab:resonant-cases} records the checked picture.

\begin{longtable}{p{0.17\linewidth}p{0.10\linewidth}p{0.31\linewidth}p{0.15\linewidth}p{0.17\linewidth}}
\toprule
Arithmetic marker & Example & Current reading & Status & Next theorem \\
\midrule
\endfirsthead
\toprule
Arithmetic marker & Example & Current reading & Status & Next theorem \\
\midrule
\endhead
$m\equiv 9,15\pmod{24}$, $5\nmid m$ & $39,57,63$ & bare sigma-core succeeds & \status{Checked} & no extra repair \\
$m\equiv 3\pmod{24}$, $5\nmid m$ & $51$ & direct reduction to width-1 pure color-$1$ frontier & \status{Checked} & pure color-$1$ breaker \\
$m\equiv 21\pmod{24}$, $5\nmid m$ & $69$ & staircase plus centered support enlargement to width-3 pure color-$1$ frontier & \status{Checked} & pure color-$1$ breaker \\
$m\equiv 21\pmod{24}$, $5\mid m$ & $45$ & $B$-active wall-localized near-miss & \status{Checked} & gate theorem \\
$m\equiv 3\pmod{24}$, $5\mid m$ & $75$ & hybrid branch: color-$0$ gain plus color-$1$ residue & \status{Checked} & hybrid gate theorem \\
\bottomrule
\caption{Current resonant case table.}
\label{tab:resonant-cases}
\end{longtable}

The table should be read honestly.
The arithmetic markers are strongly supported by exact checked packets, but the
only theorem-level statements so far are the quotient/visibility laws and the
obstruction packets described below.

\section{Pure color-\texorpdfstring{$1$}{1} frontier}

The pure color-$1$ frontier is the main place where the resonant theorem is
currently concentrated.
The checked $51$-type and $69$-type branches both reduce to canonical packages
that are already simple enough to analyze by exact section returns.

\subsection{Canonical width-1 and width-3 packages}

The two canonical pure color-$1$ packages are:
\[
\wdone=
\{A_{1,4},\,C_{r,4},\,A_{1,3},\,C_{3,3},\,A_{r,3}\},
\]
\[
\wdthree=
\{A_{1,4},\,C_{r,4},\,A_{1,3},\,C_{3,3},\,C_{r-1,3},\,A_{r,3},\,C_{r+1,3}\},
\qquad r=(m-1)/2.
\]
Color-$1$ quotient reduction shows that the type-$0$ collar rows are represented
canonically by $C$ rather than $\rho^{-1}$.
Thus, for color $1$, these really are the only canonical width-$1$ and width-$3$
models one needs to test.

\subsection{Checked deterministic reductions from \texorpdfstring{$51$}{51} and \texorpdfstring{$69$}{69}}

The current exact reductions are:
\begin{enumerate}[label=\textup{(\arabic*)},leftmargin=2.8em]
\item \textbf{$51$-type.}
      Starting from the sigma-core
      \[
      K_m=\{A_{1,4},\,C_{r,4},\,\rho_{1,3},\,C_{3,3}\},
      \]
      the checked sequence
      \[
      K_m\to K_m+A_{r,3}\to K_m+A_{r,3}+B_{1,3}
      \]
      leads to a package with exact profile
      \[
      [51^5,\ 342077451,\ 51^5]
      \]
      at $m=51$.
      This is the width-$1$ pure color-$1$ normal form.
\item \textbf{$69$-type.}
      The checked sequence
      \[
      K_m\to C_{r-1,3}\to B_{3,3}\to A_{r+1,3}\to C_{r+1,3}\to A_{r,3}
      \]
      leads to a package with exact profile
      \[
      [69^5,\ 548867331,\ 69^5]
      \]
      at $m=69$.
      This is the width-$3$ pure color-$1$ normal form.
\end{enumerate}
These reductions are \status{checked exact packets}, not all-$m$ theorems.
Their value is that they show the pure color-$1$ frontier is not an ad hoc
family zoo: it already has canonical width-$1$ and width-$3$ models.

\subsection{Width-1 as an explicit obstruction model}

The first theorem-level fact is negative.

\begin{theorem}[width-1 top-row obstruction]
Assume $3\mid m$ and consider the color-$1$ section return
\[
R^{\mathrm{wd1}}_{1,m}=h_1^m\big|_{\Sigma=0}
\]
for the package $\wdone$.
For each $e\in\{0,\dots,m-2\}$ the sets
\[
S_{0,e}=\{(0,b,c,m-1,e):c\equiv 0\pmod3\},
\qquad
S_{2,e}=\{(0,b,c,m-1,e):c\equiv 2\pmod3\}
\]
are invariant and each carries a cycle of length $m/3$.
Hence the width-$1$ pure color-$1$ frontier is not Hamilton.
\end{theorem}

\begin{proof}[Idea]
On these sets the exact width-$1$ section return collapses to
\[
(0,b,c,m-1,e)\mapsto (0,b-3,c+3,m-1,e),
\]
so one sees an explicit top-row $+3$ resonance.
This is why width $1$ should be viewed as an obstruction model rather than as a
candidate solution.
\end{proof}

\subsection{Width-3 as a sparse collar correction}

The next theorem-level fact explains why width $3$ is the natural next model.
It is not a new species of dynamics.
It is width $1$ plus a sparse correction.

\begin{theorem}[exact width-3 collar correction]
For the color-$1$ section returns of $\wdone$ and $\wdthree$ one has
\[
R^{\mathrm{wd3}}_{1,m}=R^{\mathrm{wd1}}_{1,m}+2\chi(e_0-e_1),
\qquad
\beta=\mathbf 1[c=m-2],\qquad
\chi=\mathbf 1[d+\beta\equiv m-1].
\]
Thus width $3$ differs from width $1$ only by rerouting two $x_1$-steps into
$x_0$-steps on the sparse collar blocks where $\chi=1$.
\end{theorem}

\begin{proof}[Idea]
The only extra rows in width $3$ are $C_{r-1,3}$ and $C_{r+1,3}$.
For color $1$ they are type-$0$ collar rows.
They do not create a new bulk map; they only change where two collar steps are
sent.
\end{proof}

A direct corollary is the top-collar transducer.
On
\[
d=m-1,\qquad e\neq m-1,
\]
writing $w=a+c$, the interior dynamics is
\[
(w,c)\mapsto \bigl(w+5,\ c+3-\mathbf 1[w=1]\bigr).
\]
This explains the first role of the factor $5$.
If $5\nmid m$, then $w\mapsto w+5$ sweeps all residues and every uninterrupted
top-collar episode eventually exits.
If $5\mid m$, unswept fibers remain.
This is why the $B$-active branch is tied so strongly to divisibility by $5$.

\subsection{The double-top obstruction}

The transducer viewpoint also makes the next obstruction explicit.
On the double-top plane
\[
d=e=m-1
\]
one gets a different sweep:
\[
(w,c)\mapsto \bigl(w+6,\ c+3-\mathbf 1[w=0]\bigr).
\]
Because every resonant modulus satisfies $3\mid m$, this leaves mod-$3$ fibers
visible.

\begin{theorem}[width-3 double-top obstruction]
For every resonant modulus $3\mid m$, the width-$3$ package $\wdthree$ has
invariant fibers
\[
\mathcal D_{\lambda,\mu}
=
\{(a,b,c,m-1,m-1): a+c\equiv \lambda\!\!\pmod3,\ c\equiv \mu\!\!\pmod3\},
\]
with $\lambda\in\{1,2\}$ and $\mu\in\{0,2\}$.
Each carries a cycle of length $m/3$.
Therefore width $3$ alone is not Hamilton either.
\end{theorem}

\begin{proof}[Idea]
The $+6$ sweep preserves $w\bmod 3$.
On the stated fibers no slip occurs, so the dynamics reduces again to a rigid
$+3$ drift in the $c$-coordinate.
\end{proof}

\subsection{What the future breaker theorem must do}

The pure color-$1$ frontier has therefore been compressed as far as possible by
bulk/collar analysis.
The remaining theorem must be a \emph{double-top breaker}.
At a minimum it must do one of the following on the plane $d=e=m-1$:
\begin{enumerate}[label=\textup{(\arabic*)},leftmargin=2.8em]
\item change $d$ or $e$ so that the double-top plane is no longer invariant;
\item move $c$ into the exit class $c\equiv 1\pmod3$;
\item create a new slip line on the classes that remain invisible to the
      canonical width-$3$ transducer.
\end{enumerate}
This is the cleanest formulation of the current pure color-$1$ frontier.
The open problem is not ``analyze width $3$ more'' but rather ``find and prove a
uniform double-top breaker''.

\section{The \texorpdfstring{$B$}{B}-active / gate branch}

The second resonant family does not reduce to a pure color-$1$ problem.
Its hallmark is that the $B$-channel has genuine color-$0$ leverage.
This family currently contains the checked $45$- and $75$-type branches.

\subsection{Why this branch is separate}

The top-collar transducer already explains the first split.
If $5\mid m$, then the $+5$ sweep on the top collar leaves unswept fibers.
Thus width-$3$ collar thickening cannot be the whole story.
A branch remains in which the next useful channel is no longer the
color-$1$-visible repair used in the pure color-$1$ frontier, but rather the
$B$-channel acting on color $0$.

\paragraph{The checked $45$-type picture.}
A same-support $B$-booster localizes color $0$ onto a thin wall and produces a
wall-localized near-miss.
The remaining question is a gate-destination theorem: how to redirect the wall
without reopening the closed colors.

\paragraph{The checked $75$-type picture.}
A $B$-active move improves color $0$ but leaves a genuine color-$1$ residue.
This is why the $75$-type branch should be viewed as hybrid rather than as a
copy of the $45$-type case.

At present these are still \status{checked exact packets} rather than theorem-level
all-$m$ statements.
The right theorem target is therefore not a single resonant repair theorem but a
separate $B$-active / gate theorem.

\section{Current theorem status and next theorem order}

Table~\ref{tab:status-ledger} records the present theorem status in the terms
most useful for further writing and proof search.

\begin{longtable}{p{0.33\linewidth}p{0.16\linewidth}p{0.41\linewidth}}
\toprule
Object & Status & Meaning \\
\midrule
\endfirsthead
\toprule
Object & Status & Meaning \\
\midrule
\endhead
Front-end slice $T0$--$T4$ & \status{Closed theorem} & Already part of the manuscript spine \\
Color-$4$ branch for $\SelStar$ & \status{Closed theorem} & Hamilton for all $m>2$ \\
Color-$3$ Route-E branch for $\SelStar$ & \status{Closed theorem} & Hamilton for all $m>2$ \\
Small odd residual packages $m=5,7,9$ & \status{Closed theorem} & Explicit packages close the residual colors \\
Nonresonant residual theorem for $\Pm$ & \status{Closed theorem} & Full closure for all odd $m\ge 11$ with $3\nmid m$ \\
Fixed-color quotient / visibility laws & \status{Closed theorem} & Search channels are determined by visible colors \\
$51$-type reduction to width-$1$ frontier & \status{Checked exact packet} & Pure color-$1$ normal form reached \\
$69$-type reduction to width-$3$ frontier & \status{Checked exact packet} & Same-support staircase plus centered support enlargement \\
Width-$1$ top-row obstruction & \status{Closed theorem} & Width $1$ is not Hamilton \\
Width-$3$ exact collar correction & \status{Closed theorem} & Width $3$ equals width $1$ plus sparse collar rerouting \\
Width-$3$ double-top obstruction & \status{Closed theorem} & Width $3$ is still not Hamilton by itself \\
$45$/$75$ branch picture & \status{Checked exact packet} & $B$-active / gate family remains separate \\
Resonant residual theorem & \status{Frontier} & Only remaining load-bearing input for the main theorem \\
\bottomrule
\caption{Current status ledger.}
\label{tab:status-ledger}
\end{longtable}

Given this ledger, the most natural theorem order is:
\begin{enumerate}[label=\textup{(\arabic*)},leftmargin=2.8em]
\item prove a uniform double-top breaker theorem for the pure color-$1$
      frontier;
\item prove a uniform $B$-active / gate theorem for the $45$/$75$-type branch;
\item combine these into the resonant residual theorem;
\item invoke the already rewritten final upgrade to obtain the unconditional
      odd-$m$, $d=5$ Hamilton decomposition theorem.
\end{enumerate}

\begin{question}[most useful next local theorem]
Among the remaining tasks, the most conceptually central one is the following:
can one find a local resonant surgery which is visible on the double-top plane
for the canonical width-$3$ pure color-$1$ frontier and provably destroys the
mod-$3$ invariant fibers without reopening colors $0$ and $2$?
A positive answer would collapse the pure color-$1$ frontier to a theorem and
leave only the separate $B$-active / gate family.
\end{question}

\section*{Conclusion}

The point of the rewrite is not merely that progress has been made.
It is that the remaining problem is now sharply shaped.
Another researcher reading the updated paper with the present note should see
exactly where the honest boundary lies, exactly which packets are already
stable, and exactly which two theorem families remain before the global odd-$m$,
$d=5$ theorem becomes unconditional.

\end{document}
